\section{A New Scaling Dimension}

The dominant story of AI progress is model scaling. Better data, algorithms,
and compute produce more capable models. This remains the primary source of
machine intelligence, and Fusion is designed to benefit from every advance it
creates.

But once multiple capable models exist, a second question becomes unavoidable:
how effectively can a system turn the available supply of intelligence into
useful work?

We call progress in the models themselves \emph{supply-side intelligence
improvement}. We call progress in understanding and satisfying the task's
changing needs \emph{demand-side allocation intelligence}. The two are
complementary:

\begin{equation}
  \text{system progress}
  \quad\longleftarrow\quad
  \underbrace{\text{model progress}}_{\text{supply}}
  \quad\times\quad
  \underbrace{\text{allocation progress}}_{\text{demand}}.
\end{equation}

The expression is not a proposed power law. It captures a strategic shift.
Better models improve the ingredients of the system. Better allocation changes
what the system can make from the same ingredients. A new model can strengthen
Fusion immediately; a better capability map, task representation, or stopping
policy can improve the value extracted from every model already in the pool.

\subsection{Capability without proportional frontier compute}

Inference-time compute is valuable when it changes the quality of a decision.
It is waste when it reproduces evidence the system already has, reasons around
an observation only the environment can provide, or continues after the result
has become stable. A system that cannot tell the difference lacks an important
form of intelligence about its own work.

For this reason, compute efficiency in Fusion is not a purchasing optimization
added after quality. It is a consequence of task understanding. The same
allocation decision determines both whether the system gains a useful
perspective and whether a model call earns its cost. The relevant outcome is a
better performance--compute frontier: greater capability without proportional
growth in scarce frontier computation.

This framing also explains why heterogeneity matters. Different models need
not be interchangeable substitutes. They can be different instruments. A
specialist can contribute a domain-specific observation; an efficient model can
search a broad space; an independent critic can find a counterexample; a
frontier model can integrate evidence and exercise judgment. The purpose is not
to prefer small models or large models. It is to assign each unit of work to
the capability with the highest decision value.

\subsection{Scaling through learning}

Allocation intelligence begins as system design, but it becomes a learning
problem. Every task offers evidence about which model helped, under which role,
given which state, at what cost, and with what downstream outcome. Over time,
Fusion can replace static assumptions about model capability with an empirical
map of conditional contribution. It can learn not only \emph{which} model to
call, but \emph{why}, \emph{when}, and \emph{whether} to call one at all.

\principle{Fusion aims to make system capability compound faster than
frontier-compute consumption.}
